\documentclass[11pt]{article}
\usepackage[a4paper,margin=23mm]{geometry}
\usepackage{amsmath,amssymb,amsthm,mathtools}
\newtheorem{theorem}{Theorem}[section]
\newtheorem{lemma}[theorem]{Lemma}
\newtheorem{corollary}[theorem]{Corollary}
\newtheorem{remark}[theorem]{Remark}
\title{GST Original-World Source Audit: Exact Calibration Gap and Actual Divisor-Supported Primitive Chow Birth}
\date{23 October 2026}
\begin{document}\maketitle
\begin{abstract}
The owner's exact question is why the corrected alpha-dependent
original-X smooth determinant construction has not completed
the Hodge conjecture. We audit its FIRST open geometric arrow,
remove a superfluous mass condition, and prove a precise
variational obstruction with an attained minimum. Then, using
a DIFFERENT actual geometric input from the prescribed target
alpha, we construct genuine original-X Chow representatives
in every codimension-two rational Hodge direction of geometric
coniveau at least one, by a support-Gysin Hodge lift and an
ACTUAL algebraic divisor on a resolution of the support.
The first result closes the inaccurate suggestion that
topological source birth plus positive energy dissipation
forces a holomorphic cycle; the second is an independently
executed target-sensitive, canonical-Betti source constructor
whose hypotheses do not presuppose a Chow cycle for alpha.
No additional axiom or abstract Hodge-to-Chow operator is used.
\end{abstract}
\section{Exact input/output contract}
Fix ANY smooth connected projective complex variety X
of dimension d, a codimension 1<=p<=d, an actual integral
K\"ahler class h represented by a smooth K\"ahler form
omega, and a SPECIFIED rational (p,p) Hodge class
alpha in H^{2p}(X,Q). The ultimate required output
is a true cycle Z_alpha in CH^p(X)_Q
with canonical Betti class alpha.

The 2026-10-23 previous GST signed-determinant construction
correctly produces, from ANY rational EVEN cohomology alpha
whether Hodge or not, actual closed integral REAL
determinantal homology cycles T_+,T_- on ORIGINAL X
and rational coefficients yielding PD(T_+-T_-)=alpha
after scaling. Thus it does NOT supply a property
distinguishing Hodge classes from all rational
cohomology classes. The remaining analytic-to-algebraic
step was a calibration-defect limit theorem.

\section{The exact and attained calibrated-cycle variational gap}
Put k=d-p and let
\[
 \varphi_k=\frac{\omega^k}{k!}.
\]
This is the Wirtinger calibration for positively oriented
complex k-planes. Write \(\mathcal I_{2k}(X)\) for the
closed integral rectifiable 2k-currents in X
(boundary zero), with integral homology class \([T]\).
Fix an integral homology class
\(\gamma\in H_{2k}(X,\mathbb Z)\) in the
image of the cycle map from closed integral currents.
Define its precise global spatial defect
\[
 \boxed{\mathfrak d_\omega(\gamma)=
 \inf_{\substack{T\in\mathcal I_{2k}(X)\\[T]=\gamma}}
  \Bigl(\mathbf M(T)-T(\varphi_k)\Bigr).
 }                                                       \tag{1}
\]
Because \(\varphi_k\) is closed, its evaluation
is a homological constant, say
\(V_\omega(\gamma)=\langle\gamma,[\varphi_k]\rangle\).

\begin{theorem}[The INFIMUM IS ATTAINED on every integral original-X class]
For every such integral gamma, there exists
a genuine closed integral current \(T_{\min}\)
in homology class gamma with
\[
 \boxed{\mathfrak d_\omega(\gamma)=
   \mathbf M(T_{\min})-V_\omega(\gamma)\ge0.}
                                                               \tag{2}
\]
In particular for any sequence in gamma,
\[
 \Delta_\omega(T_j)\to0
 \quad\Longrightarrow\quad
 \mathbf M(T_j)\to V_\omega(\gamma),
                                                               \tag{3}
\]
so a SEPARATE uniform-mass hypothesis is REDUNDANT.
\end{theorem}
\begin{proof}
Choose a minimizing sequence of closed integral
currents in the fixed class gamma. Because
\(\mathbf M(T_j)=\Delta_\omega(T_j)
   +V_\omega(\gamma)\) and the defects
are eventually bounded by definition of infimum,
their masses are uniformly bounded.
X is compact and their boundaries have
zero mass, so Federer--Fleming integral-current
compactness yields an integral flat limit
\(T_{\min}\). Integral homology class is
constant under flat convergence on compact X:
sufficiently small integral flat differences
represent zero integral homology class (the
realization of integral-current homology
as singular integral homology).
Therefore \([T_{\min}]=\gamma\).
Lower semicontinuity of mass and continuity
of evaluation on the smooth closed form
\(\varphi_k\) show the limit realizes the
infimum. Wirtinger calibration gives its
nonnegativity. Finally (3) follows from
the constant evaluation equality, independent
of any compactness theorem.
\end{proof}

\begin{theorem}[Zero calibration gap IFF a POSITIVE original-X algebraic cycle represents gamma]
For any integral homology class gamma as above,
\[
 \boxed{
 \mathfrak d_\omega(\gamma)=0
 \quad\Longleftrightarrow\quad
 \gamma=[C]_{\rm hom}
 \text{ for some effective INTEGRAL
 algebraic codimension-p cycle }C\text{ on X}.
 }                                                       \tag{4}
\]
The implication from zero gap to a cycle
constructs an ACTUAL \(C\) on ORIGINAL X,
not on an auxiliary folded variety.
\end{theorem}
\begin{proof}
If the defect vanishes, take the attained
minimizer T_min in (2).
Equality in the pointwise Wirtinger inequality
says that T_min is a positively
complex-calibrated, closed integral current.
The positive closed rectifiable integral-current
structure theorem (King / Harvey--Shiffman
holomorphic chain theorem) identifies it with
a positive INTEGRAL holomorphic k-cycle.
Because X is projective and the cycle is closed,
Chow's theorem algebraizes it as an effective
integral Chow cycle C with the same homology class.
Conversely, the integration current of any
effective integral complex algebraic cycle C
is pointwise calibrated and has defect exactly
zero, so its class gamma has gap zero.
\end{proof}

\begin{remark}[The exact limitation of the preceding signed topological construction]
Given a nonzero primitive rational (p,p) class r,
its pairing with \(h^{d-p}\) is zero.
Therefore any scaled INTEGRAL homology
class gamma_r with PD proportional to r
cannot itself have zero calibration gap
unless gamma_r=0: an effective nonzero
analytic cycle has positive ample degree.
For the two background-shifted
classes gamma_plus, gamma_minus in the
previous signed construction,
zero gap for EACH implies a genuine
signed algebraic realization of alpha.
The previous construction proves
gamma_plus-minus is alpha;
it DOES NOT prove both gaps vanish.
Equation (4) makes the missing
geometric statement mathematically exact,
and blocks promoting 'Hodge type' alone
to a vanishing variational certificate.
This is a corrected diagnosis,
not a claim that the entire GST approach fails.
\end{remark}

\section{ONE-GAP total effective-completion theorem: a truly geometric criterion for all X,p}
The previous signed-two-cycle theorem used
separate calibrated limits for two
topological bundles. In fact, allowing
the genuine original-X algebraic
complete-intersection background
removes ONE of those two
calibration existence requirements.

Choose the given very ample integral
divisor H with c1(H)=h, so h^p
is represented by an actual effective
complete-intersection cycle
for every p. Let alpha be ANY rational
even-degree class of degree 2p.
For integers m>0 and n>=0, an INTEGRAL LIFT of
m alpha+n h^p means an ACTUAL class
\(\gamma\in H_{2(d-p)}(X,\mathbb Z)\)
whose rational Poincare dual is m alpha+n h^p.
The lift must be specified in integral homology,
not as a canonical integral representative of
rational cohomology modulo torsion.

\begin{theorem}[ONE-GAP algebraic cycle recovery, all original varieties]
For EVERY original X, p and rational
alpha of degree 2p, the following
are equivalent, with NO hypothesis
that alpha is Hodge:
\[
 \boxed{
 \begin{split}
 &\exists\,Z_\alpha\in CH^p(X)_{\mathbb Q},
   \quad\operatorname{cl}^{can}_X(Z_\alpha)=\alpha\\
 &\qquad\Longleftrightarrow\
 \exists\,m\in\mathbb Z_{>0},\
 n\in\mathbb Z_{\ge0},\
 \gamma\in H_{2(d-p)}(X,\mathbb Z):
 \quad\operatorname{PD}_{\mathbb Q}(\gamma)=m\alpha+n h^p,
 \quad\mathfrak d_\omega(\gamma)=0.
 \end{split}
 }                                                       \tag{8}
\]
When the right-hand condition holds,
the attained zero-gap current constructs
a genuine EFFECTIVE integral projective
cycle C with
\(\operatorname{cl}^{can}_X(C)=m\alpha+n h^p\).
Any actual complete intersection
\(B\in CH^p(X)\) with
\(\operatorname{cl}^{can}_X(B)=h^p\)
then gives the entirely explicit
original-X rational Chow output
\[
 \boxed{Z_\alpha=\frac1m(C-nB).}               \tag{9}
\]
\end{theorem}
\begin{proof}
For the reverse implication use
the attained calibrated current
theorem (4) and subtract the
explicit genuine algebraic
complete intersection nB.
For the forward implication choose
a genuine rational Chow class Z
with cl(Z)=alpha and clear a
positive integer m so
mZ=A-B with A and B effective
INTEGRAL Chow cycles. Write
B as a finite sum of integral
codimension-p subvarieties
with positive integer coefficients.

For each integral component
\(W\subset X\) of B, Serre
global generation of its ideal
sheaf in a sufficiently high
power of H constructs p
global hypersurface sections
cutting out a pure
codimension-p COMPLETE
INTERSECTION \(Q_W\)
containing W with generic
intersection multiplicity one.
Consequently
\[
 [Q_W]=r_W^pH^p=[W]+[R_W]
       \quad\text{in actual }CH^p(X),
\]
where R_W is an EFFECTIVE
cycle. Multiplying each by
the multiplicity in B and
adding gives
\[
 nH^p=B+R
 \quad\text{in }CH^p(X),
 \qquad n=\sum_W
       (\operatorname{mult}_B W)r_W^p,
\]
with R effective and n a
nonnegative integer.
Therefore
\(mZ+nH^p=A+R\)
is an EFFECTIVE
algebraic integral cycle.
Its Betti class is
m alpha+n h^p, and the
associated integral positive
complex current is calibrated,
so its gap vanishes.
This proves the direct implication,
using a genuine original-X
algebraic hypersurface
excess-component construction
rather than an assumed universal
Hodge-to-Chow map.
\end{proof}

\begin{corollary}[Exact consequence for the owner's proposed infinite-dimensional experiment]
For rational (p,p) input alpha, a
UNIFORM geometric GST theorem forcing
one *attained* zero-gap class
m alpha+n h^p, with m>0 and n>=0
and an explicit actual ample H,
would construct the original-X
Hodge Chow representative (9)
in one shot. There is NO need
for an independently bounded
mass assumption and no need
to realize a second
topological bundle as a
calibrated cycle: the
complete-intersection
background is already algebraic.

Conversely, equation (8)
also shows that proving this
zero-gap existence for EVERY
rational Hodge alpha is
mathematically as strong as
the original all-X
algebraicity target. The
variational identity is
an exact translation to
spatial geometry, not
a derivation of the
universal zero-gap law.
\end{corollary}

\section{Different alpha-dependent route: codimension TWO divisor-support geometry}
The variational audit prevents repeating the
same open geometric implication under a new
functional. Now return to the original X
and use a NEW target-sensitive geometric input:
actual algebraic SUPPORT of the prescribed
Hodge class alpha.

Let p=2 and X have arbitrary complex dimension d>=2.
Suppose there exist finitely many smooth projective
varieties \(Y_i\) of dimension d-1 and
proper morphisms \(f_i:Y_i\to X\)
whose images are contained in algebraic
divisors of X. These Y_i can be chosen as
resolutions of the irreducible components
of an actual algebraic divisor support.
Assume the specified original-X
rational class alpha belongs to the
ACTUAL sum of rational Gysin images
\[
 \boxed{
 \alpha\in\sum_i
 (f_i)_*H^2(Y_i,\mathbb Q)
       \subset H^4(X,\mathbb Q).
 }                                                       \tag{5}
\]
This is a condition about algebraic SUPPORT
and maps of COHOMOLOGY, not a pre-existing
codimension-two Chow representative of alpha.
It is true, for example, whenever alpha is
a coniveau-one class supported through such a
divisor resolution.

\begin{theorem}[NEW actual original-X codimension-two Chow source from arbitrary divisor support]
For EVERY original smooth projective complex X,
EVERY rational (2,2) Hodge alpha satisfying
the independently testable support relation (5),
there are rational divisor classes
\(\beta_i\in H^2(Y_i,\mathbb Q)\cap H^{1,1}(Y_i)\)
and actual rational ALGEBRAIC divisors
\(D_i\in CH^1(Y_i)_\mathbb Q\) with
\[
 \boxed{
 \alpha=\sum_i (f_i)_*\beta_i,\qquad
 Z_\alpha=\sum_i (f_i)_*[D_i]
  \in CH^2(X)_\mathbb Q,\qquad
 \mathrm{cl}^{can}_X(Z_\alpha)=\alpha .
 }                                                       \tag{6}
\]
All images use actual proper algebraic Gysin
and Chow pushforward. No original-X Chow
preimage of alpha occurs in the hypotheses.
The theorem holds for any d, without a
restriction to special CM, K3 or elliptic worlds.
\end{theorem}
\begin{proof}
Put \(V=\bigoplus_i H^2(Y_i,\mathbb Q)(-1)\).
The map \(F=\sum_i(f_i)_*:V\to H^4(X,\mathbb Q)\)
is a morphism of POLARIZABLE pure rational
Hodge structures of weight four.
Because f_i is a proper morphism from
smooth projective dimension d-1
to smooth projective dimension d,
its Gysin raises degree by two and
bidegree by (1,1).
The image \(W=F(V)\) is a rational
Hodge substructure. The category of
polarizable pure rational Hodge structures
is semisimple: there is a rational
Hodge-equivariant section of the
surjection \(V\to W\).
Apply it to the specified rational
(2,2) alpha in W. It supplies a
rational \(\beta=(\beta_i)\) of Hodge
type (2,2) IN V, equivalently each
\(\beta_i\) has type (1,1) in
\(H^2(Y_i,\mathbb Q)\), and
\(\sum_i(f_i)_*\beta_i=\alpha\).
The Lefschetz (1,1) theorem
and projective GAGA construct an
ACTUAL algebraic rational Cartier
divisor \(D_i\) on each smooth
projective Y_i with canonical Betti
class \(\beta_i\).
Push those actual divisor cycles
along the proper f_i to obtain
the genuine codimension-two
rational Chow cycle \(Z_\alpha\)
on ORIGINAL X. The proper
pushforward compatibility of
canonical cycle classes then
gives exactly (6).
\end{proof}

\begin{corollary}[Prescribed primitive coniveau-one original-X geometry]
If 2p=4<=d, alpha is a NONZERO primitive
rational (2,2) class, and (5) holds,
the actual original-X cycle in (6) has
nonzero native primitive pairing
\[
 \boxed{
 \int_X\mathrm{cl}^{can}(Z_\alpha)\smile
            \alpha\smile h^{d-4}
 =\int_X\alpha^2 h^{d-4}>0.
 }                                                       \tag{7}
\]
The strict sign is the Hodge--Riemann
positivity for p=2.
Unlike the smooth-real topological
determinantal carrier, this source is
GENUINELY algebraic before applying the
original-X native projection or
Chow-return compiler.
\end{corollary}

\section{Sharp input boundary and actual new attempt}
The geometric operation in (6) is a genuine
target-dependent Chow SOURCE, not an
existing-source Gram inversion. But it
requires an algebraic divisor support
on which alpha admits the actual
degree-two Gysin lift (5).
An arbitrary rational primitive
(2,2) Hodge alpha may have no
known such support. The general Hodge
condition alone does not supply (5),
and replacing that with an assumption
of its universal validity would be an
unjustified source-birth postulate.
The constructive route (6) closes
the source problem for the entire
divisor-supported sector, including
any primitive directions that occur
there; it leaves the completely
unresolved original-X residual OUTSIDE
the proven Gysin image.

The owner-directed remaining
geometric research should attempt
to generate a genuine divisor-support
Gysin representation of a prescribed
residual alpha from actual GST
causal geometry, rather than
retesting fixed PGL tangent Chern
sources or creating another
equivalent 'existence of Chow'
criterion.
\end{document}